%STEPBEGIN{0}
\documentclass{standalone}
\usepackage{hsbcommon}
\begin{document}
\showmytreeraw[4]{forest-work-breakdown.tex}   %%% <-- EDIT step list; keep this filename in sync with the file's own name
\end{document}
%STEPEND
% Step 1: tasks only
% Step 2: + persons responsible for each work package (2nd line of the WP node)
% Step 3: + assignee per task
% Step 4: + duration per task
% The WP person lines sit INSIDE the node's {...} label, so filtering them
% out never touches the tree's bracket structure.
    \begin{forest}
        for tree={
            font=\footnotesize,
            grow = east,
            anchor=west,
            align=left,
            forked edge,
            inner sep=2pt
                }
        [Project
            [{Work package ...
                %STEPBEGIN{2,3,4}
                \\ Peter
                %STEPEND
             }
                [...][...]
            ]
            [{Work package 3
                %STEPBEGIN{2,3,4}
                \\ Olivia
                %STEPEND
             }
                %STEPBEGIN{1,2}
                [Task 3.2 ]
                [Task 3.1 ]
                %STEPEND
                %STEPBEGIN{3}
                [Task 3.2 - Alice]
                [Task 3.1 - Alice]
                %STEPEND
                %STEPBEGIN{4}
                [Task 3.2 - Alice - 12d]
                [Task 3.1 - \textit{Alice - 15d}]
                %STEPEND
            ]
            [{Work package 2
                %STEPBEGIN{2,3,4}
                \\ Mona
                %STEPEND
             }
                %STEPBEGIN{1,2}
                [Task 2.2 ]
                [Task 2.1 ]
                %STEPEND
                %STEPBEGIN{3}
                [Task 2.2 - Charlie]
                [Task 2.1 - Bob]
                %STEPEND
                %STEPBEGIN{4}
                [Task 2.2 - Charlie - 6d]
                [Task 2.1 - Bob - 8d]
                %STEPEND
            ]
            [{Work package 1
                %STEPBEGIN{2,3,4}
                \\ Lisa
                %STEPEND
             }
                %STEPBEGIN{1,2}
                [Task 1.3 ]
                [Task 1.2 ]
                [Task 1.1 ]
                %STEPEND
                %STEPBEGIN{3}
                [Task 1.3 - Doreen]
                [Task 1.2 - Alice]
                [Task 1.1 - Doreen]
                %STEPEND
                %STEPBEGIN{4}
                [Task 1.3 - Doreen - 9d]
                [Task 1.2 - Alice - 4d]
                [Task 1.1 - Doreen - 9d]
                %STEPEND
            ]
        ]
    \end{forest}